% ch04_u03 —— 原书 p.142–p.151（4.2.11 末、4.2.12–4.2.16、4.3.1–4.3.2）
% 首句承接 ch04_u02 末的 a) 项（p.141 末），故本节自 b) 项起；用 start=2 接续原书编号。
\SourcePage{142}
\begin{enumerate}[label=\alph*),leftmargin=2.2em,start=2]
\item (179) 与来自 4.2.8 的配对
$\operatorname{Pf}(V;L_{+},L_{-})^{*}\otimes\operatorname{Pf}(V^{-};L_{+},L_{-})^{*}\xrightarrow{\ \sim\ }\det(V;L_{+},L_{-})^{*}$
及
$\operatorname{Pf}(\tilde{V};\tilde{L}_{+},\tilde{L}_{-})^{*}\otimes\operatorname{Pf}(\tilde{V}^{-};L_{+},L_{-})^{*}\xrightarrow{\ \sim\ }\det(\tilde{V};\tilde{L}_{+},\tilde{L}_{-})^{*}$
相容。
% [Tutor Note: 原书 p.142 第二个配对中，Pfaffian 自变量里的 L 缺波浪号（第一个配对中是带波浪号的 L），底本如此，保留原样。]
\end{enumerate}

b) 可直接验证，a) 由 b) 推出。也可以把 $(\tilde{V};\tilde{L}_{+},\tilde{L}_{-})$ 归到它是\emph{极大子商}的情形（即 $\Lambda\oplus(L_{+}\cap L_{-})=L_{+}$），再用 4.2.5(ii) 证明 a)。

\bdp{4.2.12} 设 $E$ 是 $X$ 上的向量丛。则 $E\oplus(E^{*}\otimes\omega_{X})$ 具有显然的 $\omega$-正交丛结构。我们将构造典范同构
\begin{equation}\tag{180}
\operatorname{Pf}(E\oplus(E^{*}\otimes\omega_{X}))\xrightarrow{\ \sim\ }\det R\Gamma(X,E).
\end{equation}
选取子层 $E'\subset E$ 与局部自由层 $E''\supset E$，使得 $H^{0}(X,E')=0$、$H^{1}(X,E'')=0$，且 $E''/E'$ 具有限支撑。置
\[
V:=H^{0}(X,E''/E'),\qquad L_{+}:=H^{0}(X,E'/E)\subset V,\qquad L_{-}:=H^{0}(X,E'')\subset V.
\]
则 $R\Gamma(X,E)$ 可等同于复形 $0\to L_{-}\to V/L_{+}\to 0$，而 $\det R\Gamma(X,E)$ 可等同于 $\det(V;L_{+},L_{-})$。另一方面，$\mathcal{Q}:=E\oplus(E^{*}\otimes\omega_{X})$ 的 Pfaffian 可用子层 $\mathcal{Q}':=E'\oplus((E'')^{*}\otimes\omega_{X})\subset\mathcal{Q}$ 来计算。于是 $\operatorname{Pf}_{\mathcal{Q}'}(\mathcal{Q})$ 等于 (170) 的左端。故 (170) 给出同构 (180)。可以验证 (180) 不依赖 $E'$ 与 $E''$。

\bdp{4.2.13} Lagrange 三元组的概念有一个有用的无穷维推广。先回忆一些基本定义。

\emph{定义.} \emph{Tate 空间}是指完备拓扑向量空间，其 $0$ 的邻域基由一族可公度的向量子空间构成（即对基中任意 $U_{1},U_{2}$，有 $\dim U_{1}/(U_{1}\cap U_{2})<\infty$）。

\emph{注.} Tate 空间隐含在 Tate 的卓越著作 [T] 之中。事实上，[T] 中发展的曲线上的留数理论，最自然的解释正是通过 Tate 空间的自同态代数\SourcePage{143}的典范中心扩张；这一点也隐含在 [T] 中。Tate 扩张的构造可参看 7.13.18。

设 $V$ 是 Tate 空间。若对每个开子空间 $U\subset V$ 都存在有限集 $\{v_{1},\dots,v_{n}\}\subset V$，使得 $P\subset U+\mathbb{C}v_{1}+\dots+\mathbb{C}v_{n}$，则称向量子空间 $P\subset V$ \emph{有界}。$V$ 的\emph{拓扑对偶}是 $V$ 上连续线性泛函构成的空间 $V^{*}$，其上的（线性）拓扑使得 $V$ 的有界子空间的正交补构成 $0\in V^{*}$ 的邻域基。显然 $V^{*}$ 是 Tate 空间，且典范态射 $V\to(V^{*})^{*}$ 是同构。

\emph{例}（坐标 Tate 空间）。设 $I$ 是集合。若 $A\setminus(A\cap B)$ 与 $B\setminus(B\cap A)$ 都有限，则称 $A,B\subset I$ \emph{可公度}。可公度性是一个等价关系。设 $\mathcal{A}$ 是子集 $A\subset I$ 的族的一个等价类。$\mathcal{A}$ 的元素称为\emph{半无限子集}。记 $\mathbb{C}((I,\mathcal{A}))$ 为形式线性组合 $\sum_{i}c_{i}e_{i}$（$c_{i}\in\mathbb{C}$）构成的空间，其中要求 $c_{i}$ 在 $i\notin A$ 时为零，而 $A$ 是某个半无限子集。它是 Tate 向量空间（拓扑由子空间 $\mathbb{C}[[A]]:=\{\sum_{i\in A}c_{i}e_{i}\}$ 定义，其中 $A$ 半无限）。$\mathbb{C}((I,\mathcal{A}))$ 的对偶空间是 $\mathbb{C}((I,\mathcal{A}'))$，其中 $\mathcal{A}'$ 由 $\mathcal{A}$ 中子集的补集构成。任何 Tate 向量空间都同构于适当选取的 $I,\mathcal{A}$ 的 $\mathbb{C}((I,\mathcal{A}))$；这样的同构由相应的子集 $\{e_{i}\}\subset V$ 给出，它称为 $V$ 的\emph{拓扑基}。

$V$ 中的 \emph{c-格}是开有界子空间。$V$ 中的 \emph{d-格}\footnote{c 与 d 分别是 “compact” 与 “discrete” 的首字母。}是离散子空间 $\Gamma\subset V$，使得对某个 c-格 $P\subset V$ 有 $\Gamma+P=V$。若 $W\subset V$ 是 d-格（相应地，c-格），则存在 c-格（相应地，d-格）$W'\subset V$，使得 $V=W\oplus W'$。若 $W\subset V$ 是 d-格（相应地，c-格），则 $W^{\perp}\subset V^{*}$ 也是 d-格（相应地，c-格），且 $(W^{\perp})^{\perp}=W$。

Tate 空间 $V$ 上的（连续）双线性型称为\emph{非退化}的，是指它诱导出拓扑同构 $V\to V^{*}$。设 $V$ 是配备非退化对称双线性型的 Tate 空间。若 $L^{\perp}=L$，则称子空间 $L\subset V$ 为 \emph{Lagrange} 的。

\SourcePage{144}
\emph{定义.} \emph{Tate Lagrange 三元组}由以下数据构成：配备非退化对称双线性型的 Tate 空间 $V$，$V$ 中的 Lagrange c-格 $L_{+}\subset V$，以及 $V$ 中的 Lagrange d-格 $L_{-}\subset V$。

\emph{例.} 设 $\mathcal{Q}$ 是 $X$ 上的 $\omega$-正交丛。对 $x\in X$，记 $\mathcal{Q}\otimes\mathcal{O}_{x}$（相应地，$\mathcal{Q}\otimes K_{x}$）为 $\mathcal{Q}$ 拉回到 $\operatorname{Spec}\mathcal{O}_{x}$（相应地，$\operatorname{Spec}K_{x}$）上之后的整体截面空间。$\mathcal{Q}\otimes K_{x}$ 是 Tate 空间，带非退化对称双线性型 $\operatorname{Res}(\ ,\ )$。对每个非空有限子集 $S\subset X$，有 Tate Lagrange 三元组
\begin{equation}\tag{181}
V:=\oplus_{x\in S}(\mathcal{Q}\otimes K_{x}),\qquad L_{+}:=\oplus_{x\in S}(\mathcal{Q}\otimes\mathcal{O}_{x}),\qquad L_{-}:=\Gamma(X\setminus S,\mathcal{Q}).
\end{equation}

设 $(V;L_{+},L_{-})$ 是 Tate Lagrange 三元组。则对任意 c-格 $\Lambda\subset L_{+}$，只要 $\Lambda\cap L_{-}=0$，就有 (178) 所定义的有限维 Lagrange 三元组 $(\tilde{V};\tilde{L}_{+},\tilde{L}_{-})$。如 4.2.11 所述，$\operatorname{Pf}(\tilde{V};\tilde{L}_{+},\tilde{L}_{-})$ 与 $\det(\tilde{V};\tilde{L}_{+},\tilde{L}_{-})$ 不依赖 $\Lambda$。置
\[
\operatorname{Pf}(V;L_{+},L_{-}):=\operatorname{Pf}(\tilde{V};\tilde{L}_{+},\tilde{L}_{-}),\qquad
\det(V;L_{+},L_{-}):=\det(\tilde{V};\tilde{L}_{+},\tilde{L}_{-}).
\]
等价地，可以把 $\det(V;L_{+},L_{-})$ 定义为复形 (164) 的行列式，而把 $\operatorname{Pf}(V;L_{+},L_{-})$ 用 (166) 或 (167) 来定义（此时应设 (166) 中的 $\operatorname{Cl}(V)$-模 $M$ 是\emph{离散}的，即对每个 $m\in M$，$\{v\in V\mid vm=0\}$ 是开集）。

\emph{例.} 若 $(V;L_{+},L_{-})$ 由 (181) 定义，则 $\operatorname{Pf}(V;L_{+},L_{-})=\operatorname{Pf}(\mathcal{Q})$，$\det(V;L_{+},L_{-})=\det R\Gamma(X,\mathcal{Q})$。

4.2.7 与 4.2.8 中的构造在 Tate 情形下仍有意义，只需作如下显然的改动：a) 在 4.2.7 中应设 $L_{+}$ 是 c-格，$L_{-}$ 是 d-格；b) (171) 应换成如下公式：
\begin{equation}\tag{182}
M=\varinjlim_{U}\bigwedge(V/U)\otimes\det(L_{+}/U)^{*}
\end{equation}
其中 $U$ 取遍 $L_{+}$ 中的 c-格。(182) 的右端是\emph{费米 Fock 空间}，即 $V$ 的\emph{半无限外幂}的直和（参看 [KR] 第 4 讲及其中的文献）。

\SourcePage{145}
\emph{注.} 用三元组 (181) 表述 $\operatorname{Pf}(\mathcal{Q})$ 的式子可以改写如下。对 $x\in X$，考虑交换李超代数 $\mathfrak{a}_{\mathcal{O}_{x}}\subset\mathfrak{a}_{K_{x}}$，使 $\mathfrak{a}_{\mathcal{O}_{x}}$（相应地，$\mathfrak{a}_{K_{x}}$）的奇部为 $\mathcal{Q}\otimes\mathcal{O}_{x}$（相应地，$\mathcal{Q}\otimes K_{x}$），偶部为 $0$。$\mathcal{Q}\otimes K_{x}$ 上的对称双线性型定义了中心扩张 $0\to\mathbb{C}\to\tilde{\mathfrak{a}}_{K_{x}}\to\mathfrak{a}_{K_{x}}\to 0$，它在 $\mathfrak{a}_{\mathcal{O}_{x}}$ 上有典范分裂。Clifford 代数 $\operatorname{Cl}(\mathcal{Q}\otimes K_{x})$ 就是扭泛包络代数 $\operatorname{U}'\mathfrak{a}_{K_{x}}$，而
\[
M_{x}:=\operatorname{Cl}(\mathcal{Q}\otimes K_{x})/\operatorname{Cl}(\mathcal{Q}\otimes K_{x})\cdot(\mathcal{Q}\otimes\mathcal{O}_{x})
\]
是 $\operatorname{U}'\mathfrak{a}_{K_{x}}$ 上的真空模。按 (167)，$\operatorname{Pf}(\mathcal{Q})^{*}$ 是 $\Gamma(X\setminus S,\mathcal{Q})$ 在 $\otimes_{x\in S}M_{x}$ 上作用的余不变量空间。

\bdp{4.2.14} 本小节讨论 Tate Lagrange 三元组的\emph{族}。设 $R$ 是交换环。\emph{Tate $R$-模}是指同构于 $P\oplus Q^{*}$ 的拓扑 $R$-模，其中 $P,Q$ 是有限生成投射 $R$-模的（无限）直和（$0\in P\oplus Q^{*}$ 的邻域基由所有有限生成子模 $M\subset Q$ 对应的 $M^{\perp}\subset Q^{*}$ 构成）。这个\emph{坏}\footnote{有限秩投射 $R((t))$-模不必是上述意义的 Tate 模。我们的 Tate $R$-模概念对 $\operatorname{Spec}R$ 不是局部的。它还有其它缺点。}定义对我们的目的够用。事实上，我们主要处理同构于 $V_{0}\hat{\otimes}R$ 的 Tate $R$-模，其中 $V_{0}$ 是 Tate 空间。

若如下定义 c-格与 d-格，则 4.2.13 中 Tate 线性代数的讨论对 Tate $R$-模仍成立。

\emph{定义.} Tate $R$-模 $V$ 中的 \emph{c-格}是开有界子模 $P\subset V$，使 $V/P$ 投射。$V$ 中的 \emph{d-格}是子模 $\Gamma\subset V$，使得对 $V$ 中的某个 c-格 $P$ 有 $\Gamma\cap P=0$，且 $V/(\Gamma+P)$ 是有限型投射模。\footnote{于是这对所有 c-格 $P'\subset P$ 也成立。}

现若 $1/2\in R$，则可完全像 4.2.13 那样定义 Tate Lagrange 三元组（当然，若 $1/2\notin R$，则应改用二次型代替双线性型，这很容易）。Tate Lagrange 三元组的 Pfaffian\SourcePage{146}$(V;L_{+},L_{-})$ 在 $R$ 上按 4.2.13 中同样的方式定义，只作如下小的修改：为过渡到 (178) 所定义的有限维 Lagrange 三元组 $(\tilde{V};\tilde{L}_{+},\tilde{L}_{-})$，须设 $\Lambda\subset L_{+}$ 是 c-格，使 $\Lambda\cap L_{-}=0$ 且 $V/(\Lambda+L_{-})$ 投射（这两个性质等价于 $\Lambda^{\perp}+L_{-}=V$）。

\emph{例.} 设 $D\subset X\otimes R$ 是 $\operatorname{Spec}R$ 上的有限闭子概形，可局部由一个方程定义（即 $D$ 是有效相对 Cartier 除子）。设 $\mathcal{Q}$ 是 $X\otimes R$ 上的向量丛。设态射 $D\to\operatorname{Spec}R$ 满射。则
\[
V:=\varprojlim_{m}\varinjlim_{n}H^{0}(X\otimes R,\mathcal{Q}(nD)/\mathcal{Q}(-mD))
\]
是 Tate $R$-模\footnote{事实上 $V$ 同构于某个 Tate 空间 $V_{0}$ 的 $V_{0}\hat{\otimes}R$。确实，可以设 $R$ 在 $\mathbb{C}$ 上有限生成，再用 7.12.11。我们需要 7.12.11 的下述情形：$R$ 在 $\mathbb{C}$ 上有限生成，且 7.12.11 中的投射模是有限生成模的直和；此时 7.12.11 由 Serre 定理（[Se] 定理 1；另见 [Ba68] 第 4 章 §2）与 Eilenberg 引理 [Ba63] 推出。}，
\[
L_{+}:=\varprojlim_{m}H^{0}(X\otimes R,\mathcal{Q}/\mathcal{Q}(-mD))\subset V
\]
是 c-格，且
\[
L_{-}:=H^{0}((X\otimes R)\setminus D,\mathcal{Q})\subset V
\]
是 d-格。若 $\mathcal{Q}$ 是 $\omega$-正交丛，则 $(V;L_{+},L_{-})$ 是 Lagrange 三元组，且 $\operatorname{Pf}(\mathcal{Q})=\operatorname{Pf}(V;L_{+},L_{-})$（参看 4.2.13）。

\bdp{4.2.15} 记 $\mathcal{B}$ 为 $\mathbb{C}$ 上配备非退化对称双线性型的有限维向量空间构成的群胚。本小节（读者可以跳过）构造典范同构
\begin{equation}\tag{183}
\begin{aligned}
&\operatorname{Pf}(V\otimes W;L_{+}\otimes W,L_{-}\otimes W)\\
&\qquad\xrightarrow{\ \sim\ }\operatorname{Pf}(V;L_{+},L_{-})^{\otimes\dim W}\otimes|\det W|^{\otimes p(V;L_{+},L_{-})},
\end{aligned}
\end{equation}
% [Tutor Note: 原书 p.146 把编号 (183) 单独印在公式上一行的左端，公式本身未带编号；底本如此，此处按 equation 排版。]
\begin{equation}\tag{184}
\operatorname{Pf}(\mathcal{Q}\otimes W)\xrightarrow{\ \sim\ }\operatorname{Pf}(\mathcal{Q})^{\otimes\dim W}\otimes|\det W|^{\otimes p(\mathcal{Q})}
\end{equation}

\SourcePage{147}
其中 $W\in\mathcal{B}$，$(V;L_{+},L_{-})$ 是（Tate）Lagrange 三元组，$\mathcal{Q}$ 是 $X$ 上的 $\omega$-正交丛，$|\det W|$ 是把 $W$ 看作空间（不是超空间！）时的行列式，而 $p(V;L_{+},L_{-}),p(\mathcal{Q})\in\mathbb{Z}/2\mathbb{Z}$ 分别是 $\operatorname{Pf}(V;L_{+},L_{-})$ 与 $\operatorname{Pf}(\mathcal{Q})$ 的奇偶性。$|\det W|^{\otimes n}$ 对 $n\in\mathbb{Z}/2\mathbb{Z}$ 有意义，因为有典范同构 $|\det W|^{\otimes 2}\xrightarrow{\ \sim\ }\mathbb{C}$，$(w_{1}\wedge\dots\wedge w_{r})^{\otimes 2}\mapsto\det(w_{i},w_{j})$。

为定义 (183) 与 (184)，注意 $\mathcal{B}$ 是以 $\oplus$ 为张量「积」的张量范畴。(183) 与 (184) 的两端都是从 $\mathcal{B}$ 到一维超空间范畴的张量函子（为把 (184) 的右端看作张量函子，把它改写为 $|\operatorname{Pf}(\mathcal{Q})|^{\otimes\dim W}\otimes(\det W)^{\otimes p(\mathcal{Q})}$，其中 $|\operatorname{Pf}(\mathcal{Q})|$ 由 $\operatorname{Pf}(\mathcal{Q})$ 改变 $(\mathbb{Z}/2\mathbb{Z})$-分次使其为偶而得到，而 $\det W$ 是把 $W$ 看作\emph{超空间}时的行列式）。

断言：把 (183) 与 (184) 定义为张量函子的同构，存在唯一的方式使得对 $W=(\mathbb{C},1)$，(183) 与 (184) 等于 $\operatorname{id}$。这里 $1$ 表示双线性型 $(x,y)\mapsto xy$，$x,y\in\mathbb{C}$。

为证明，对由 (183) 或 (184) 的左端除以其右端得到的张量函子 $F$，应用下述引理。

\emph{引理.} 每个张量函子 $F:\mathcal{B}\to\{1\text{ 维向量空间}\}$ 都同构于张量函子 $F_{1}$，后者由 $F_{1}(W)=L^{\otimes\dim W}$ 定义，其中 $L:=F(\mathbb{C},1)$。存在唯一同构 $F\xrightarrow{\ \sim\ }F_{1}$，它诱导恒等映射 $F(\mathbb{C},1)\to F_{1}(\mathbb{C},1)$。

\emph{证明.} 对每个 $W\in\mathcal{B}$，函子 $F$ 诱导同态 $f_{W}:\operatorname{Aut}W\to\mathbb{C}^{*}$。因 $\operatorname{Aut}W$ 是正交群，故 $f_{W}(g)=(\det g)^{n(W)}$，其中 $n(W)\in\mathbb{Z}/2\mathbb{Z}$。显然 $n(W)=n$ 与 $W$ 无关。置 $W_{1}:=(\mathbb{C},1)$。$F$ 把交换性同构 $\left(\begin{smallmatrix}0&1\\1&0\end{smallmatrix}\right):W_{1}\oplus W_{1}\to W_{1}\oplus W_{1}$ 映为 $\operatorname{id}$。故 $n=0$，即 $f_{W}$ 对每个 $W$ 平凡。其余是显然的，因为 $\mathcal{B}$ 的对象同构类的半群 $|\mathcal{B}|$ 就是 $\mathbb{Z}_{+}$。\hfill$\Box$

\emph{注.}
\begin{enumerate}[label=(\roman*),leftmargin=2.2em]
\item (183) 在 4.2.8 中已被隐含地使用。
\item 我们将在 4.2.16 中使用 (183)。
\end{enumerate}

\SourcePage{148}
\bdp{4.2.16} 本小节（读者当然可以跳过）解释把 $\mathbb{C}$ 换成特征 $2$ 的域 $k$ 时会发生什么。此时必须区分二次型（见 [Bourb59] §3, n°4）与对称双线性型。在 Lagrange 三元组的定义中，$V$ 应带非退化\emph{二次}型。于是 $\omega$-正交丛的定义中，$\mathcal{Q}$ 应带非退化\emph{二次}型 $\mathcal{Q}\to\omega_{X}$（因 $k$ 的特征为 $2$，非退化性蕴含 $\mathcal{Q}$ 的秩为偶）。4.2.8 中 $\operatorname{Pf}\otimes\operatorname{Pf}^{-}\xrightarrow{\ \sim\ }\det$ 的构造须修改。若 $(V;L_{+},L_{-})$ 是 Lagrange 三元组，而 $W$ 带非退化对称\emph{双线性}型，则 $(V\otimes W;L_{+}\otimes W,L_{-}\otimes W)$ 是 Lagrange 三元组。双线性型 $\left(\begin{smallmatrix}1&0\\0&-1\end{smallmatrix}\right)$ 与 $\left(\begin{smallmatrix}0&1\\1&0\end{smallmatrix}\right)$ 在特征 $2$ 下不等价，但可以对 $W=H$ 与 $W=H'$ 使用 (183) 来构造 $\operatorname{Pf}\otimes\operatorname{Pf}^{-}\xrightarrow{\ \sim\ }\det$。最后，须在特征 $2$ 下构造 (183) 与 (184)。为简单起见设 $k$ 完全。则 (183) 与 (184) 的同构的特征性质\footnote{在非完全情形表述此性质时，应把 $\mathcal{B}$ 看作栈而非群胚。}的表述与 4.2.15 相同，但存在性与唯一性的证明须修改。半群 $|\mathcal{B}|$（见 4.2.15 引理证明的末尾）不再是 $\mathbb{Z}_{+}$；它有生成元 $a,b$ 与定义关系 $a+b=3a$（$a$ 对应 $1$ 阶矩阵 $(1)$，$b$ 对应 $\left(\begin{smallmatrix}0&1\\1&0\end{smallmatrix}\right)$）。故 $\mathcal{B}$ 对应的群是 $\mathbb{Z}$，这已够用。

\bdsec{4.3. Pfaffian（二）}

\bdp{4.3.1} 固定 $\mathbb{C}$ 上 $n$ 维向量空间 $W$ 与其上的非退化对称双线性型 $(\ )$。为简化记号，把 $O(W)$ 与 $SO(W)$ 写作 $O_{n}$ 与 $SO_{n}$。

设 $\mathcal{F}$ 是 $X$ 上的 $SO_{n}$-torsor。相应的秩 $n$ 向量丛 $W_{\mathcal{F}}$ 带双线性型 $(\ )_{\mathcal{F}}$，且有典范同构 $\det W_{\mathcal{F}}=\mathcal{O}_{X}\otimes\det W$。设 $\mathcal{L}\in\omega^{1/2}(X)$，即 $\mathcal{L}$ 是 $\omega_{X}$ 的平方根。

\SourcePage{149}
则 $W_{\mathcal{F}}\otimes\mathcal{L}$ 是 $\omega$-正交丛，故 $\operatorname{Pf}(W_{\mathcal{F}}\otimes\mathcal{L})$ 有意义（见 4.2）。考虑「归一化」Pfaffian
\begin{equation}\tag{185}
\operatorname{Pf}_{\mathcal{L},\mathcal{F}}:=\operatorname{Pf}(W_{\mathcal{F}}\otimes\mathcal{L})\otimes\operatorname{Pf}(W\otimes\mathcal{L})^{\otimes-1}
\end{equation}
与「归一化」行列式
\begin{equation}\tag{186}
\nu(\mathcal{F}):=\det R\Gamma(X,W_{\mathcal{F}})\otimes\det R\Gamma(X,\mathcal{O}_{X}\otimes W)^{\otimes-1}.
\end{equation}
如 4.2.1 所述，存在典范同构 $c_{\pm i}:\operatorname{Pf}^{\otimes 2}\xrightarrow{\ \sim\ }\det R\Gamma$。使用例如 $c_{i}$ 得到同构\footnote{故同构 (187)=(189) 依赖 $-1$ 的平方根的选取。若把 (187) 乘以 $i^{\pm p(\mathcal{F})^{2}}$，则此依赖消失，其中 $p$ 是典范映射 $\operatorname{Bun}_{SO_{n}}\to\pi_{0}(\operatorname{Bun}_{SO_{n}})=\pi_{1}(SO_{n})=\mathbb{Z}/2\mathbb{Z}$，而 $p(\mathcal{F})^{2}\in\mathbb{Z}/4\mathbb{Z}$。由于 4.2.1 的注 (iii) 所述的理由，我们宁愿不这样做。}
\begin{equation}\tag{187}
\operatorname{Pf}_{\mathcal{L},\mathcal{F}}^{\otimes 2}\xrightarrow{\ \sim\ }\nu_{\mathcal{L}}(\mathcal{F})
\end{equation}
其中
\begin{equation}\tag{188}
\nu_{\mathcal{L}}(\mathcal{F}):=\det R\Gamma(X,W_{\mathcal{F}}\otimes\mathcal{L})\otimes\det R\Gamma(X,W\otimes\mathcal{L})^{\otimes-1}.
\end{equation}
[Del87] 中的构造 7.2 给出典范同构
\[
\nu_{\mathcal{L}}(\mathcal{F})=\nu(\mathcal{F})\otimes\langle\det W_{\mathcal{F}}\otimes(\det W)^{\otimes-1},\mathcal{L}\rangle
\]
因 $\det W_{\mathcal{F}}=\mathcal{O}_{X}\otimes\det W$，故 $\nu_{\mathcal{L}}(\mathcal{F})=\nu(\mathcal{F})$，且
\begin{equation}\tag{189}
\operatorname{Pf}_{\mathcal{L},\mathcal{F}}^{\otimes 2}=\nu(\mathcal{F}).
\end{equation}
$\mathcal{F}$ 变化时，$\operatorname{Pf}_{\mathcal{L},\mathcal{F}}$ 与 $\nu(\mathcal{F})$ 是 $\operatorname{Bun}_{SO_{n}}$ 上两个线丛的纤维，二者分别记作 $\operatorname{Pf}_{\mathcal{L}}$ 与 $\nu$。

记 $\nu^{1/2}(\operatorname{Bun}_{SO_{n}})$ 为 $\nu$ 的平方根构成的范畴。有函子
\begin{equation}\tag{190}
\operatorname{Pf}:\omega^{1/2}(X)\to\nu^{1/2}(\operatorname{Bun}_{SO_{n}})
\end{equation}
它由 $\mathcal{L}\mapsto\operatorname{Pf}_{\mathcal{L}}$ 定义。

\SourcePage{150}
$\omega^{1/2}(X)$ 与 $\nu^{1/2}(\operatorname{Bun}_{SO_{n}})$ 是 Picard 范畴 $\mu_{2}\operatorname{tors}(X)$ 与 $\mu_{2}\operatorname{tors}(\operatorname{Bun}_{SO_{n}})$ 上的 torsor。有 Picard 函子 $\ell^{\operatorname{Spin}}:\mu_{2}\operatorname{tors}(X)\to\mu_{2}\operatorname{tors}(\operatorname{Bun}_{SO_{n}})$；这就是 4.1 中的函子 $\ell=\ell^{\tilde{G}}$ 在 $G=SO_{n}$、$\tilde{G}=\operatorname{Spin}_{n}$、$\Pi=\mathbb{Z}/2\mathbb{Z}$ 时的特例。在 4.3.8–4.3.15 中我们将证明，函子 $\operatorname{Pf}:\omega^{1/2}(X)\to\nu^{1/2}(\operatorname{Bun}_{SO_{n}})$ 具有典范的 $\ell^{\operatorname{Spin}}$-仿射函子结构。在此之先，4.3.2–4.3.7 将证明：对有限子集 $S\subset X$，4.1.7 所定义的 $SO_{n}(K_{S})$ 在 $\operatorname{Bun}_{SO_{n},S}$ 上的作用，提升为 $SO_{n}(K_{S})$ 的某个中心扩张在 $\operatorname{Pf}_{\mathcal{L}}$ 到 $\operatorname{Bun}_{SO_{n},S}$ 的拉回上的作用。一旦引入此作用，就很容易用 $SO_{n}(K_{S})$-不变性来刻画函子 $\operatorname{Pf}$ 上的 $\ell^{\operatorname{Spin}}$-仿射结构（见 4.3.8–4.3.10）。

\bdp{4.3.2} 设 $V$ 是带\emph{偶型}非退化对称双线性型的 Tate 空间，即存在 Lagrange c-格 $L\subset V$（见 4.2.13）；若 $\dim V<\infty$，这表示 $\dim V$ 为偶。记 $O(V)$ 为 $V$ 的保持该型的拓扑自同构群。回忆典范中心扩张的熟知构造
\begin{equation}\tag{191}
0\to\mathbb{C}^{*}\to\tilde{O}(V)\to O(V)\to 0.
\end{equation}
设 $M$ 是 Clifford 代数 $\operatorname{Cl}(V)$ 上的不可约 $(\mathbb{Z}/2\mathbb{Z})$-分次离散模（离散是指：对每个 $m\in M$，$\{v\in V\mid vm=0\}$ 是开集）。则 $M$ 在张量一维 $(\mathbb{Z}/2\mathbb{Z})$-分次空间的意义下唯一。故 $O(V)$ 在 $M$ 中有自然的射影表示。(191) 就是与此表示对应的扩张，即
\[
\tilde{O}(V):=\{(g,\varphi)\mid g\in O(V),\ \varphi\in\operatorname{Aut}_{\mathbb{C}}M,\ \varphi(vm)=g(v)\cdot\varphi(m)\ \text{对}\ m\in M\}.
\]
显然 $\tilde{O}(V)$ 不依赖 $M$ 的选取（事实上 $\operatorname{Aut}_{\mathbb{C}}M$ 是 $\operatorname{Cl}(V)$ 的自然完备化的可逆元群）。若 $(g,\varphi)\in\tilde{O}(V)$，则 $\varphi$ 或为偶或为奇。记 $\chi(g)\in\mathbb{Z}/2\mathbb{Z}$ 为 $\varphi$ 的奇偶性。则 $\chi:O(V)\to\mathbb{Z}/2\mathbb{Z}$ 是同态。

\SourcePage{151}
$-1\in O(V)$ 在 $\tilde{O}(V)$ 中的原像不是中心的。事实上，若 $\varphi:M\to M$ 满足：$m$ 为偶时 $\varphi(m)=m$，$m$ 为奇时 $\varphi(m)=-m$，则 $[-1]:=(-1,\varphi)\in\tilde{O}(V)$，且
\begin{equation}\tag{192}
[-1]\cdot\tilde{g}=(-1)^{\chi(g)}\cdot\tilde{g}\cdot[-1],\qquad g\in O(V)
\end{equation}
其中 $\tilde{g}$ 表示 $g$ 在 $\tilde{O}(V)$ 中的一个原像。

$O(V)$ 与 $\operatorname{Aut}_{\mathbb{C}}M$ 有自然的群归纳概形结构。更确切地说，把 $\mathbb{C}$-代数 $R$ 对应到集合 $O(V\hat{\otimes}R)$ 与 $\operatorname{Aut}_{\mathbb{C}}(M\hat{\otimes}R)$ 的函子是归纳概形（若 $\dim V=\infty$，则它们可表示为一族\emph{不可数}的滤过闭子概形之并）。故 $\tilde{O}(V)$ 是群归纳概形。

记 $\operatorname{Lagr}(V)$ 为 $V$ 中 Lagrange c-格构成的集合。它有自然的归纳概形结构：$\operatorname{Lagr}(V)=\varinjlim_{\Lambda}\operatorname{Lagr}(\Lambda^{\perp}/\Lambda)$，其中 $\Lambda$ 取遍 $V$ 中的迷向 c-格（故 $\operatorname{Lagr}(V)$ 的 $R$-点就是 4.2.14 意义下 $V\hat{\otimes}R$ 中的 Lagrange c-格）。记 $\mathcal{P}=\mathcal{P}_{M}$ 为 $\operatorname{Lagr}(V)$ 上的线丛，它在 $L\in\operatorname{Lagr}(V)$ 上的纤维等于 $M^{L}:=\{m\in M\mid Lm=0\}$。$O(V)$ 在 $\operatorname{Lagr}(V)$ 上的作用典范地提升为 $\tilde{O}(V)$ 在 $\mathcal{P}$ 上的作用。

$\operatorname{Lagr}(V)$ 有两个连通分支，由一维 $(\mathbb{Z}/2\mathbb{Z})$-分次空间 $M^{L}$（$L\in\operatorname{Lagr}(V)$）的奇偶性区分。此断言的证明容易化归到 $\dim V$ 有限（且为偶）的情形。同样的论证表明：$L_{1},L_{2}\in\operatorname{Lagr}(V)$ 属于同一分支，当且仅当 $\dim(L_{1}/(L_{1}\cap L_{2}))$ 为偶。显然，$\operatorname{Lagr}(V)$ 的连通分支对 $g\in O(V)$ 不变，当且仅当 $\chi(g)=0$。因此 $\chi:O(V)\to\mathbb{Z}/2\mathbb{Z}$ 是群归纳概形的态射。

我们来证明 (191) 来自群归纳概形的正合列
\begin{equation}\tag{193}
0\to\mathbb{G}_{m}\to\tilde{O}(V)\to O(V)\to 0.
\end{equation}
只需说明态射 $\tilde{O}(V)\to O(V)$ 是 $\mathbb{G}_{m}$-torsor。为此固定 $L\in\operatorname{Lagr}(V)$，并置 $M=\operatorname{Cl}(V)/\operatorname{Cl}(V)L$，于是 $\mathcal{P}=\mathcal{P}_{M}$ 在 $L$ 上的纤维等于 $\mathbb{C}$。定义 $f:O(V)\to\operatorname{Lagr}(V)$ 为 $f(g)=gL$。
